\SourceLeaf{311}
\[
\begin{aligned}
\log b={}&\log\beta+\log n+\tfrac18nn\lambda\lambda+\tfrac1{64}n^4\lambda^4+\tfrac1{384}n^6\lambda^6+\tfrac1{2048}n^8\lambda^8+\text{u.\,s.\,w.}\\
&+\tfrac1{24}\beta\beta(nn+\tfrac14n^4\lambda\lambda+\tfrac1{16}n^6\lambda^4+\tfrac1{64}n^8\lambda^6+\text{u.\,s.\,w.})\\
&+\tfrac{11}{2880}\beta^4(n^4+\tfrac12n^6\lambda\lambda+\tfrac3{16}n^8\lambda^4+\text{u.\,s.\,w.})\\
&+\tfrac{191}{362880}\beta^6(n^6+\tfrac34n^8\lambda\lambda+\text{u.\,s.\,w.})\\
&+\tfrac{2497}{29030400}\beta^8(n^8+\text{u.\,s.\,w.})\\
&+\text{u.\,s.\,w.}
\end{aligned}
\]
oder, indem man die Glieder gleicher Ordnung zusammenfasst,
\[
\begin{aligned}
\log b={}&\log\beta+\log n\\
&+\tfrac1{24}nn(\beta\beta+3\lambda\lambda)\\
&+\tfrac1{2880}n^4(11\beta^4+30\beta\beta\lambda\lambda+45\lambda^4)\\
&+\tfrac1{362880}n^6(191\beta^6+693\beta^4\lambda\lambda+945\beta\beta\lambda^4+945\lambda^6)\\
&+\tfrac1{29030400}n^8(2497\beta^8+11460\beta^6\lambda\lambda+20790\beta^4\lambda^4\\
&\qquad\qquad+18900\beta\beta\lambda^6+14175\lambda^8)\\
&+\text{u.\,s.\,w.}
\end{aligned}
\]

\indent Um die Gleichungen I, II, III in eine ganz entwickelte Gestalt zu bringen, wird man in denselben noch substituiren
\[
\begin{aligned}
\log m&=\tfrac1{12}rr+\tfrac7{1440}r^4+\tfrac{31}{90720}r^6+\tfrac{127}{4838400}r^8+\text{u.\,s.\,w.}\\
mm&=1+\tfrac16rr+\tfrac{17}{720}r^4+\tfrac{31}{10080}r^6+\text{u.\,s.\,w.}\\
m^4&=1+\tfrac13rr+\tfrac3{40}r^4+\text{u.\,s.\,w.}\\
m^6&=1+\tfrac12rr+\text{u.\,s.\,w.}\\
&\text{u.\,s.\,w.}\\
\log n&=-\tfrac1{24}rr-\tfrac1{2880}r^4-\tfrac1{181440}r^6-\tfrac1{9676800}r^8-\text{u.\,s.\,w.}\\
nn&=1-\tfrac1{12}rr+\tfrac1{360}r^4-\tfrac1{20160}r^6+\text{u.\,s.\,w.}\\
n^4&=1-\tfrac16rr+\tfrac1{80}r^4-\text{u.\,s.\,w.}\\
n^6&=1-\tfrac14rr+\text{u.\,s.\,w.}\\
&\text{u.\,s.\,w.}
\end{aligned}
\]

\indent Wir erhalten demnach für die Logarithmen von \(t,l,b\), oder vielmehr für die Unterschiede dieser Logarithmen von den genäherten Werthen \(\log\tau,\log\lambda,\log\beta\), zusammengesetzte Reihen; welche fortschreiten
\begin{list}{}{\setlength{\leftmargin}{1.4em}\setlength{\itemindent}{0pt}\setlength{\itemsep}{0pt}\setlength{\parsep}{0pt}}
\item für \(\log t\) nach den geraden Potenzen von \(\tau\) und \(r\), und deren Producten,
\item für \(\log l\) eben so nach \(\lambda\) und \(r\),
\item für \(\log b\) nach \(\beta,\lambda\) und \(r\),
\end{list}
